\documentclass[10pt,a4paper]{amsart}
\usepackage[margin=19mm]{geometry}
\usepackage{amsmath,amssymb,mathtools}
\usepackage[scr=rsfs,cal=euler]{mathalfa}
\usepackage{xfrac}
\usepackage[unicode,hidelinks,bookmarks=false]{hyperref}
\newcommand{\ie}{i.e.\ }
\newcommand{\cf}{cf.\ }
\newcommand{\resp}{respectively\ }
\newcommand{\Subsection}{\subsection}
\providecommand{\nobreakdash}{\nobreak\hbox{-}}
\setlength{\emergencystretch}{2em}
\multlinegap0pt
\usepackage{amssymb}
\usepackage{mathtools}
\usepackage{float}
\usepackage{xcolor}
\usepackage[shortlabels]{enumitem}
\usepackage{algorithm}\usepackage{algpseudocode}
\usepackage{tcolorbox}
\newtheorem{proposition}{Proposition}
\title{Securing the Flow: Maritime Energy Resilience under Correlated and Decision-Dependent Disruptions}
\author{Monit Sharma, Hoong Chuin Lau}
\date{Source: arXiv:2605.11990v1 (CC BY 4.0)}
\begin{document}
\maketitle

\noindent\emph{Source and licence.} Securing the Flow: Maritime Energy Resilience under Correlated and Decision-Dependent Disruptions. Version arXiv:2605.11990v1 (CC BY 4.0), distributed by arXiv under the Creative Commons Attribution 4.0 International licence, http://creativecommons.org/licenses/by/4.0/, as recorded in the arXiv licence metadata for this exact version and shown on the abstract page https://arxiv.org/abs/2605.11990v1. Full licence notice: \texttt{licenses/arxiv-2605.11990-CC-BY-4.0.txt}.\par\smallskip

\noindent\textit{Editorial scope.} Counted unit: the article's proposition prop:cvar-decomp (source0.tex lines 921--928). The article's complete original proof of it is delivered with it (source0.tex lines 930--967, one \texttt{proof} environment of 38 lines). No mathematics is written, repaired or re-derived: the statement and the proof are the article's own lines, in the article's own notation, and no retained line is substituted or re-flowed.  No further source-local context is needed: every cross-reference of every retained line resolves inside the two delivered spans.  Nothing was omitted from any retained span, so the package carries no omission record.  No retained line contains a citation, so the wrapper carries no bibliography and no bibliography entry.  No retained line contains a figure environment, an image-inclusion command or drawing code, so no image is delivered and the package declares no asset dependency.  Editorial formatting only, in addition: an AMS \texttt{amsart} preamble that restates the article's own declarations and macros used by the retained lines, namely the environments \texttt{proposition} on separate counters, together with the standard packages those lines need.  The retained environments print in this document as Proposition 1 in order of appearance, whereas the article prints the same items with the numbering of its own full document.  The complete native source of the article as distributed in the arXiv e-print for this exact version is bundled unchanged as \texttt{sources/200396/source0.tex}.
\begin{proposition}[Scenario-wise decomposability under mean--CVaR]
\label{prop:cvar-decomp}
For fixed first-stage decisions $(\mathbf{y}, \mathbf{W}, \nu)$, the
second-stage problem remains separable across scenarios. In particular,
for each scenario $s \in S$, the associated recourse variables and CVaR
excess variable $\xi_s$ can be optimized independently of all other
scenarios.
\end{proposition}

\begin{proof}
Fix $(\mathbf{y}, \mathbf{W}, \nu)$. Then the decision-dependent
probabilities $p_s(\mathbf{y})$ are constants, and the first-stage
variables no longer appear as optimization variables in the recourse
stage.

The uncertain part of the objective is
\[
(1-\lambda)\sum_{s\in S} p_s(\mathbf{y}) Q_s
\;+\;
\frac{\lambda}{1-\alpha}\sum_{s\in S} p_s(\mathbf{y}) \xi_s,
\]
since the term $\lambda \nu$ is constant once $\nu$ is fixed. For each
scenario $s$, the CVaR constraints are
\[
\xi_s \ge Q_s - \nu, \qquad \xi_s \ge 0.
\]
Hence, at optimality,
\[
\xi_s = \max\{Q_s-\nu,0\},
\]
so the contribution of scenario $s$ to the objective is
\[
(1-\lambda)p_s(\mathbf{y})Q_s
\;+\;
\frac{\lambda}{1-\alpha}p_s(\mathbf{y})\max\{Q_s-\nu,0\}.
\]

Moreover, the recourse constraints for scenario $s$ (flow balance,
capacity, admissibility, storage, and corridor-cap constraints) involve
only scenario-$s$ variables. No recourse variable from scenario $s$ appears
in the constraints of any other scenario $s' \neq s$.

Therefore, once $(\mathbf{y}, \mathbf{W}, \nu)$ is fixed, the second-stage
optimization decomposes into independent scenario subproblems, one for
each $s \in S$. The only cross-scenario coupling is through the shared
first-stage decisions, which are handled in the master problem.
\end{proof}

\end{document}
